% P055 简要说明（与 01--06 正文配套）
\textbf{结论名称：} P055 直线反射与斜率变换规律

\textbf{解决什么问题：} 已知一条直线（或角的边）关于另一条直线翻折，不必先求对称点坐标，就能判断翻折后直线的方向或斜率；特别适合内心、角平分线、反射及曲线对称题。

\textbf{适用条件：} 平面直角坐标系内的轴对称。先把“直线的方向变化”与“整条直线的位置变化”区分开；涉及斜率时，两条直线均须不是竖直线，否则应改用方向增量表示。

\textbf{核心快记：} 原直线斜率为 $k$，反射后的斜率为 $k'$。
\[
\boxed{\begin{array}{c|c}
\text{对称轴}&\text{斜率变化}\\ \hline
x=a\text{ 或 }y=b&k'=-k\\
y=x+c\text{ 或 }y=-x+c&k'=\dfrac1k\ (k\ne0)
\end{array}}
\]
水平线反射时“上下变号”，竖直线反射时“左右变号”；与坐标轴成 $45^\circ$ 的直线反射时“横纵交换”（可能同时变号）。遇到水平、竖直线互换时，斜率 $0$ 与“不存在”要单独处理。

\textbf{任意非竖直对称轴：} 若对称轴为 $y=mx+c$，原直线斜率为有限实数 $k$，则
\[
\boxed{k'=\frac{2m+(m^2-1)k}{1-m^2+2mk}}
\]
仅在分母非零时使用；分母为零表示反射后的直线竖直。$c$ 决定反射后的位置，不影响方向变化。

\textbf{一眼应用：} 在三角形 $PAB$ 中，若内心 $I$ 与顶点 $P$ 同在竖直直线 $x=a$ 上，则 $PI$ 是角平分线，射线 $PA,PB$ 关于 $x=a$ 对称；当斜率存在时 $k_{PA}=-k_{PB}$。

\textbf{最小提醒：} “两条斜率相反”只保证方向符合水平或竖直镜面反射，不足以单凭斜率断定它们关于\emph{指定的}轴对称；还要检查相对位置、交点及射线方向。
